\begin{trapbox}
\textbf{误区一：外角就把邻边比倒过来。}
错误写成 $BD/DC=AC/AB$，会把点放到另一侧。
下图固定同一三角形 $AB=6,AC=4,BC=5$，正确点为 $D$，错误倒置得到的点为 $D_*$。
\begin{center}
\begin{tikzpicture}[font=\small]
\draw[black!30,rounded corners=2pt] (0,0) rectangle (14.4,8.2);
\node[anchor=west,font=\small\bfseries] at (0.5,7.72) {比例倒置会改变位置：同一三角形的正确点 $D$ 与错误点 $D_*$};
% MODEL 6 4 5
\begin{scope}[shift={(5.5,4)},scale=0.43]
\coordinate (A) at (4.5,3.9686269666);
\coordinate (B) at (0,0);
\coordinate (C) at (5,0);
\coordinate (D) at (15,0);
\coordinate (W) at (-10,0);
\draw[black!45] (W)--(D);
\draw[black!45] (-0.6,0)--(15.6,0);
\draw[blue,thick] (A)--(B);
\draw[teal!80!black,thick] (A)--(C);
\draw[black,thick] (B)--(C);
\draw[purple,very thick] (A)--(D);
\fill (A) circle (1.5pt) node[above=4pt] {$A$};
\fill (B) circle (1.5pt) node[below left=4pt] {$B$};
\fill (C) circle (1.5pt) node[below right=4pt] {$C$};
\fill (D) circle (1.5pt) node[below right=4pt] {$D$};
\draw[red!75!black,thick,dashed] (A)--(W);\fill[red!75!black] (W) circle (1.5pt) node[below=5pt] {$D_*$};


\end{scope}
\node[anchor=west,text=blue] at (1,5.7) {$AB=6$};
\node[anchor=west,text=teal!80!black] at (10,5.7) {$AC=4$};
\node[anchor=west,text=red!75!black] at (0.8,6.6) {错误倒置：$10/15=2/3$};
\node[anchor=west,text=purple] at (8.3,6.6) {正确：$15/10=6/4$};
\draw[blue,|-|] ($(B)+(0,-1.1)$)--($(D)+(0,-1.1)$) node[midway,below=3pt] {$BD=15$};
\draw[teal!80!black,|-|] ($(C)+(0,-2.0)$)--($(D)+(0,-2.0)$) node[midway,below=3pt] {$DC=10$};
\node[anchor=west,text=red!75!black] at (0.8,1.2) {$AD_*$ 不是外角平分线};

\end{tikzpicture}
\end{center}

正确替代：按字母对应记 $BD\leftrightarrow AB$、$DC\leftrightarrow AC$。
当 $AB>AC$ 时必须 $BD>DC$，所以在 $C$ 外侧；
图中 $D_*$ 虽也在延长线上，但 $BD_*/D_*C=2/3$，不满足正确的 $3/2$。
仅“点在外部”不足以判定它在外角平分线上。

\textbf{误区二：把外点当成内点，用了和关系。}
错误套用 $BD+DC=BC$，相当于把例 3 中的外点 $D$ 错当成内点 $E$。
例 3 图中橙线 $AE$ 对应 $BE+EC=BC$，紫线 $AD$ 对应 $BD-DC=BC$。
正确替代：先写顺序，再写差；不确定左右时统一写 $|BD-DC|=BC$。

\textbf{误区三：只因三点共线，就说两个角互补。}
共线还不够，必须看射线方向。
$D$ 在外部时，$DB,DC$ 同向，故 $\angle ADB=\angle ADC$；
$E$ 在内部时，$EB,EC$ 反向，才有 $\angle AEB+\angle AEC=180^\circ$。
正确替代：按正弦证明图先辨同向、反向，再判断相等或互补。

\textbf{误区四：等腰仍硬求交点，或混淆有向线段与普通长度。}
$AB=AC$ 时外角平分线平行于 $BC$，没有有限 $D$；不能把分母 $|p-q|=0$ 解释为可用的长度。
若另行采用有向线段，须声明约定。例如取 $B=0,C=a$，定义
$\overline{BD}=x_D$、$\overline{DC}=a-x_D$，则
\[
 \frac{\overline{BD}}{\overline{DC}}=-\frac pq.
\]
负号来自方向，不能塞进本条以正长度书写的 $BD/DC=p/q$。
\end{trapbox}
